\subsection{Jordan and Canonical Forms}

% **Role:** structure revealed by invariant subspaces and eigenbehavior.
%
% Jordan form, rational/canonical form if desired, complexification, real block forms.

Jordan and canonical forms turn primary decomposition into explicit blocks. Jordan chains, complexification, and real block forms organize the nilpotent and real spectral structure.

\begin{definition} \label{def:jordan-chain-block-and-matrix} \hlt{Jordan Chains, Blocks, and Matrices}\\
For $k\geq1$, the \hlt{nilpotent Jordan block} $J_k(0)$ is the $k\times k$ matrix with $1$ on the first superdiagonal and $0$ elsewhere:
\begin{align}
J_k(0)=
\begin{bmatrix}
0&1&0&\cdots&0\\
0&0&1&\ddots&\vdots\\
\vdots&\ddots&\ddots&\ddots&0\\
0&\cdots&0&0&1\\
0&\cdots&\cdots&0&0
\end{bmatrix}.
\nonumber
\end{align}
For $\lambda\in\mathbb{C}$, set $J_k(\lambda)\defeq\lambda I_k+J_k(0)$. A \hlt{Jordan matrix} is a direct sum of Jordan blocks. A basis is a \hlt{Jordan basis} for an operator $T$ if the matrix of $T$ in that basis is a Jordan matrix.
\end{definition}

\begin{perspective} \label{per:jordan-strings-as-columns} \hlt{Jordan Strings as Columns}\\
A length-$k$ Jordan string for $A$ at $\lambda$ is $(\mathbf{x}_1,\ldots,\mathbf{x}_k)$ with $(A-\lambda I)\mathbf{x}_1=\mathbf{0}$; $(A-\lambda I)\mathbf{x}_j=\mathbf{x}_{j-1}$, $2\leq j\leq k$.\\
For $X=[\mathbf{x}_1\ \cdots\ \mathbf{x}_k]$, $AX=XJ_k(\lambda)$.\\
If all strings form the columns of a basis matrix $M$, then $AM=MJ$ and $M^{-1}AM=J$.
\end{perspective}

\begin{figure}[H]
  \centering
  \begin{tikzpicture}[
  >=Stealth,
  every node/.style={font=\small},
  vec/.style={draw=blue!65!black, fill=blue!7, rounded corners=2pt,
    minimum height=0.55cm, minimum width=0.95cm, align=center},
  arrow/.style={->, blue!70!black, line width=0.75pt},
  note/.style={text=gray!70!black, font=\footnotesize}
]
  \node[vec] (x1) at (0,0) {$\mathbf{x}_1$};
  \node[vec] (x2) at (1.63,0) {$\mathbf{x}_2$};
  \node[vec] (dots) at (3.21,0) {$\cdots$};
  \node[vec] (xk) at (4.80,0) {$\mathbf{x}_k$};

  \draw[arrow] (x1) -- ++(-0.90,0) node[left] {$\mathbf{0}$};
  \draw[arrow] (x2) -- node[above] {$N$} (x1);
  \draw[arrow] (dots) -- node[above] {$N$} (x2);
  \draw[arrow] (xk) -- node[above] {$N$} (dots);
\end{tikzpicture}

  \caption{Jordan chain for $N=A-\lambda I$; columns give $AX=XJ_k(\lambda)$.}
  \label{fig:jordan-string-columns-similarity}
\end{figure}

\begin{method} \label{meth:filippov-column-space-jordan-strings} \hlt{Filippov Column-Space Jordan Strings}\\
Let $A\in\operatorname{Mat}(n,n,\mathbb{C})$, let $\lambda\in\operatorname{spec}(A)$, and set $N=A-\lambda I_n$ on the generalized eigenspace $\mathcal{G}_{\lambda}(A)=\ker(N^n)$.
\begin{enumerate}[label=(\roman*)]
\setlength{\itemsep}{0pt}
\item Work first inside $\operatorname{range}(N)\cap\mathcal{G}_{\lambda}(A)$, where $N$ is again nilpotent.
\item Build Jordan strings for this smaller nilpotent operator.
\item For each string ending in $\mathbf{u}$, solve $N\mathbf{v}=\mathbf{u}$ in $\mathcal{G}_{\lambda}(A)$ and lift the string by appending $\mathbf{v}$.
\item Add enough vectors from $\ker(N)$ to complete a basis of $\mathcal{G}_{\lambda}(A)$; these added vectors are length-one strings.
\item Repeat the process recursively until all of $\mathcal{G}_{\lambda}(A)$ is covered, then concatenate the strings for all eigenvalues.
\end{enumerate}
Theorem \ref{thm:nilpotent-jordan-chain-basis} justifies the lifts; Perspective \ref{per:jordan-strings-as-columns} turns the resulting columns into $AM=MJ$.
\end{method}

\begin{proposition} \label{prop:powers-of-jordan-block} \hlt{Powers of a Jordan Block}\\
For $p\geq0$, $J_k(\lambda)^p =\sum\limits_{r=0}^{\min(p,k-1)}\binom{p}{r}\lambda^{p-r}J_k(0)^r$. Thus $J_k(\lambda)^p\to0$ if and only if $\abs{\lambda}<1$.
\begin{proof}
\begin{enumerate}[label=\roman*.]
\setlength{\itemsep}{0pt}
\item[$(\Rightarrow)$] If $J_k(\lambda)^p\to0$, then its diagonal entries $\lambda^p$ tend to $0$, so $\abs{\lambda}<1$.
\item[$(\Leftarrow)$] If $\abs{\lambda}<1$, then for each fixed $0\leq r<k$ the scalar $\binom{p}{r}\lambda^{p-r}$ tends to $0$, because a polynomial in $p$ is dominated by the exponential decay of $\abs{\lambda}^{p}$.
\end{enumerate}
\hlt{Structural Commentary}: Jordan powers are exponential-polynomial terms.\\
\hlt{Key Insight}: Polynomial growth cannot defeat exponential decay when $\abs{\lambda}<1$.
\end{proof}
\end{proposition}

\begin{theorem} \label{thm:nilpotent-jordan-chain-basis} \hlt{Nilpotent Jordan Chain Basis}\\
Let $\mathcal{V}$ be finite-dim, $N\in\LL(\mathcal{V})$ be nilpotent. Then $\exists$ vectors $\mathbf{v}_1,\ldots,\mathbf{v}_r\in\mathcal{V}$ and nonnegative integers $m_1,\ldots,m_r$ such that $\mathcal{B}_N= \bigl(N^{m_1}\mathbf{v}_1,\ldots,N\mathbf{v}_1,\mathbf{v}_1,\ldots, N^{m_r}\mathbf{v}_r,\ldots,N\mathbf{v}_r,\mathbf{v}_r\bigr)$ is a basis of $\mathcal{V}$ and $N^{m_j+1}\mathbf{v}_j=\mathbf{0}$ for every $j$.
\begin{proof}
Induct on $n=\dim(\mathcal{V})$. If $N=0$, any basis works with all $m_j=0$.\\
Now assume $N\neq0$. Then $N$ is not injective, so $\operatorname{range}(N)$ is a proper subspace, and $N|_{\operatorname{range}(N)}$ is nilpotent.\\
By induction, $\operatorname{range}(N)$ has a chain basis $\bigl(N^{m_1}\mathbf{u}_1,\ldots,N\mathbf{u}_1,\mathbf{u}_1,\ldots, N^{m_r}\mathbf{u}_r,\ldots,N\mathbf{u}_r,\mathbf{u}_r\bigr)$ with $N^{m_j+1}\mathbf{u}_j=\mathbf{0}$.\\
Choose $\mathbf{v}_j\in\mathcal{V}$ such that $N\mathbf{v}_j=\mathbf{u}_j$, and lift the chains to $\mathcal{L}= \bigl(N^{m_1+1}\mathbf{v}_1,\ldots,N\mathbf{v}_1,\mathbf{v}_1,\ldots, N^{m_r+1}\mathbf{v}_r,\ldots,N\mathbf{v}_r,\mathbf{v}_r\bigr)$.\\
The list $\mathcal{L}$ is linearly independent. Indeed, applying $N$ to a dependence among vectors in $\mathcal{L}$ gives a dependence among the range-chain basis, so all coefficients except possibly those on the leading vectors $N^{m_j+1}\mathbf{v}_j$ vanish.\\
The remaining relation is a relation among $N^{m_j+1}\mathbf{v}_j=N^{m_j}\mathbf{u}_j$, which are members of the range-chain basis, so those coefficients vanish too. Also $N(\operatorname{span}(\mathcal{L}))=\operatorname{range}(N)$: each range-chain basis vector $N^q\mathbf{u}_j$ is $N(N^q\mathbf{v}_j)$ for the corresponding lifted vector. Extend $\mathcal{L}$ to a basis by adding $\mathbf{w}_1,\ldots,\mathbf{w}_s$. For each $i$, choose $\mathbf{x}_i\in\operatorname{span}(\mathcal{L})$ such that $N\mathbf{x}_i=N\mathbf{w}_i$, and set $\mathbf{z}_i=\mathbf{w}_i-\mathbf{x}_i$. Then $N\mathbf{z}_i=\mathbf{0}$. Since each $\mathbf{x}_i$ lies in $\operatorname{span}(\mathcal{L})$, replacing $\mathbf{w}_i$ by $\mathbf{z}_i$ preserves the span and linear independence of the full basis. Thus $\mathcal{L},\mathbf{z}_1,\ldots,\mathbf{z}_s$ is the desired chain basis.\\
\hlt{Structural Commentary}: Build chains in $\operatorname{range}(N)$, lift them once, then add kernel chains.\\
\hlt{Key Insight}: Applying $N$ turns lifted-chain dependences into known range-chain dependences.
\end{proof}
\end{theorem}

\begin{figure}[H]
  \centering
  \begin{tikzpicture}[
  font=\small,
  vec/.style={draw=blue!65!black, fill=blue!7, rounded corners=2pt, minimum height=0.58cm, minimum width=1.12cm, align=center},
  arrow/.style={-{Stealth[length=2mm]}, thick, blue!70!black}
]
  \node[vec] (a0) at (0,0) {$N^m\mathbf{v}$};
  \node[vec] (a1) at (1.76,0) {$N^{m-1}\mathbf{v}$};
  \node[vec] (a2) at (3.34,0) {$\cdots$};
  \node[vec] (a3) at (4.88,0) {$N\mathbf{v}$};
  \node[vec] (a4) at (6.38,0) {$\mathbf{v}$};

  \draw[arrow] (a4) -- node[above] {$N$} (a3);
  \draw[arrow] (a3) -- node[above] {$N$} (a2);
  \draw[arrow] (a2) -- node[above] {$N$} (a1);
  \draw[arrow] (a1) -- node[above] {$N$} (a0);
  \draw[arrow] (a0) -- ++(-1.18,0) node[left] {$\mathbf{0}$};
\end{tikzpicture}

  \caption{A nilpotent chain shifts left under $N$; in this basis order it becomes one Jordan block.}
  \label{fig:jordan-chain-nilpotent-basis}
\end{figure}

\begin{theorem} \label{thm:complex-jordan-form} \hlt{Jordan Form over $\mathbb{C}$}\\
Let $\mathcal{V}\neq\{\mathbf{0}\}$ be finite-dimensional over $\mathbb{C}$ and let $T\in\LL(\mathcal{V})$. Then $\mathcal{V}$ has a Jordan basis for $T$.
\begin{proof}
Let $\lambda_1,\ldots,\lambda_m$ be the distinct eigenvalues of $T$. By Theorem \ref{thm:complex-primary-decomposition-generalized-eigenspaces}, $\mathcal{V}=\mathcal{G}_{\lambda_1}(T)\oplus\cdots\oplus \mathcal{G}_{\lambda_m}(T)$, and each $N_j\defeq(T-\lambda_jI_{\mathcal{V}})|_{\mathcal{G}_{\lambda_j}(T)}$ is nilpotent. Apply Theorem \ref{thm:nilpotent-jordan-chain-basis} to $N_j$ on each generalized eigenspace.\\
The resulting chain basis of $\mathcal{G}_{\lambda_j}(T)$ is a basis of that summand, and the direct-sum decomposition makes the concatenation of these bases a basis of $\mathcal{V}$. In a chain of length $r$ for $N_j$, the matrix of $N_j$ is $J_r(0)$, so the matrix of $T|_{\mathcal{G}_{\lambda_j}(T)}=\lambda_jI+N_j$ on the same chain is $J_r(\lambda_j)$. Therefore the concatenated basis is a Jordan basis for $T$.\\
\hlt{Structural Commentary}: Primary pieces decompose into nilpotent chains.\\
\hlt{Key Insight}: Jordan form is ``eigenvalue plus nilpotent chains'' on every generalized eigenspace.
\end{proof}
\end{theorem}

\begin{figure}[H]
  \centering
  \begin{tikzpicture}[
  font=\small,
  block/.style={draw=blue!65!black, fill=blue!6, rounded corners=2pt, minimum width=2.15cm, minimum height=0.68cm, align=center},
  arrow/.style={-{Stealth[length=2mm]}, thick, gray!75}
]
  \node[block] (v) at (0,0) {$\mathcal{V}$};
  \node[block] (g1) at (3.0,0.85) {$\mathcal{G}_{\lambda_1}(T)$};
  \node[block] (gj) at (3.0,0) {$\cdots$};
  \node[block] (gm) at (3.0,-0.85) {$\mathcal{G}_{\lambda_m}(T)$};

  \node[block] (j1) at (6.0,0.85) {$J_k(\lambda_1)$};
  \node[block] (jj) at (6.0,0) {$\cdots$};
  \node[block] (jm) at (6.0,-0.85) {$J_k(\lambda_m)$};

  \draw[arrow] (v.east) -- (g1.west);
  \draw[arrow] (v.east) -- (gj.west);
  \draw[arrow] (v.east) -- (gm.west);
  \draw[arrow] (g1.east) -- (j1.west);
  \draw[arrow] (gj.east) -- (jj.west);
  \draw[arrow] (gm.east) -- (jm.west);
\end{tikzpicture}

  \caption{Primary pieces $\mathcal{G}_{\lambda_j}(T)$ give nilpotent chains, hence Jordan blocks.}
  \label{fig:jordan-form-primary-chains}
\end{figure}

\begin{corollary} \label{cor:jordan-blocks-diagonalizability} \hlt{Diagonalizability in Jordan Form}\\
Let $A\in\operatorname{Mat}(n,n,\mathbb{C})$ have Jordan form $J$. For each eigenvalue $\lambda$, $\dim E_{\lambda}(A)$ equals the number of $\lambda$-blocks in $J$. Hence $A$ is diagonalizable if and only if every Jordan block has size $1$.
\begin{proof}
Write $A=SJS^{-1}$. Then $S:E_{\lambda}(J)\to E_{\lambda}(A)$ is an isomorphism. A block $J_k(\mu)$ contributes to $E_{\lambda}(J)$ only when $\mu=\lambda$, and then $\dim\ker(J_k(\lambda)-\lambda I_k)=\dim\ker J_k(0)=1$.\\
Thus $\dim E_{\lambda}(A)$ is the number of $\lambda$-blocks. By Theorem \ref{thm:diagonalizable-iff-algebraic-geometric-multiplicities-match}, diagonalizability means geometric and algebraic multiplicities agree for every $\lambda$; here this holds exactly when every block has size $1$.\\
\hlt{Structural Commentary}: Each Jordan block contributes one eigenvector direction.\\
\hlt{Key Insight}: Missing eigenvectors are measured by block sizes greater than $1$.
\end{proof}
\end{corollary}

\begin{lemma} \label{lem:nilpotent-jordan-block-power-rank} \hlt{Powers of a Nilpotent Jordan Block}\\
For $k\geq1$ and $p\geq0$, $\operatorname{rank}\bigl(J_k(0)^p\bigr)=\max(k-p,0)$ and $\dim\ker\bigl(J_k(0)^p\bigr)=\min(p,k)$.
\begin{proof}
Let $N=J_k(0)$. In the standard basis, $N\mathbf{e}_1=\mathbf{0}$ and $N\mathbf{e}_j=\mathbf{e}_{j-1}$ for $2\leq j\leq k$. Iterating gives
\begin{align}
N^p\mathbf{e}_j=
\begin{cases}
\mathbf{0}, & j\leq p,\\
\mathbf{e}_{j-p}, & j>p.
\end{cases} \nonumber
\end{align}
If $p<k$, the image is $\operatorname{span}(\mathbf{e}_1,\ldots,\mathbf{e}_{k-p})$, so the rank is $k-p$; if $p\geq k$, the image is $\{\mathbf{0}\}$, so the rank is $0$.\\
Hence $\operatorname{rank}(N^p)=\max(k-p,0)$. Rank-nullity gives $\dim\ker(N^p)=k-\max(k-p,0)=\min(p,k)$.\\
\hlt{Structural Commentary}: A nilpotent Jordan block is a left-shift on the standard basis.\\
\hlt{Key Insight}: Each power deletes one more leading layer of the chain.
\end{proof}
\end{lemma}

\begin{definition} \label{def:weyr-characteristic} \hlt{Weyr Characteristic}\\
Let $A\in\operatorname{Mat}(n,n,\mathbb{C})$ and $\lambda\in\operatorname{spec}(A)$.\\
For $p\geq1$, define $w_p(\lambda,A)\defeq \operatorname{rank}\bigl((A-\lambda I_n)^{p-1}\bigr) - \operatorname{rank}\bigl((A-\lambda I_n)^p\bigr)$.\\
The positive part of the sequence $w_1,w_2,\ldots$ is the \hlt{Weyr characteristic} of $A$ associated with $\lambda$.
\end{definition}

\begin{theorem} \label{thm:jordan-block-sizes-from-rank-drops} \hlt{Jordan Block Sizes from Rank Drops}\\
Let $A\in\operatorname{Mat}(n,n,\mathbb{C})$, $\lambda\in\operatorname{spec}(A)$, and $p\geq1$. Then $w_p(\lambda,A)$ is the number of $\lambda$-Jordan blocks of size at least $p$; hence $w_p(\lambda,A)-w_{p+1}(\lambda,A)$ is the number of $\lambda$-blocks of size exactly $p$.
\begin{proof}
Rank is similarity-invariant, so assume $A$ is in Jordan form.\\
Blocks $J_k(\mu)$ with $\mu\neq\lambda$ contribute $0$ to the rank drop because $J_k(\mu)-\lambda I_k$ is invertible.\\
A block $J_k(\lambda)$ contributes $\operatorname{rank}\bigl(J_k(0)^{p-1}\bigr)-\operatorname{rank}\bigl(J_k(0)^p\bigr)$, which is $1$ exactly when $p\leq k$ by Lemma \ref{lem:nilpotent-jordan-block-power-rank}.\\
Summing over blocks proves the first claim; subtracting the $p+1$ count gives the exact-size count.\\
\hlt{Structural Commentary}: Rank drops count surviving chain lengths.\\
\hlt{Key Insight}: A length-$k$ block contributes to drops $1,\ldots,k$.
\end{proof}
\end{theorem}

\begin{figure}[H]
  \centering
  \begin{tikzpicture}[
  font=\small,
  rank/.style={draw=blue!60!black, fill=blue!6, rounded corners=2pt, minimum width=0.96cm, minimum height=0.50cm, align=center},
  count/.style={draw=red!60!black, fill=red!6, rounded corners=2pt, minimum width=2.90cm, minimum height=0.50cm, align=center},
  arrow/.style={-{Stealth[length=2mm]}, thick, gray!75}
]
  \node[rank] (r0) at (0,0) {$r_0$};
  \node[rank] (r1) at (1.36,0) {$r_1$};
  \node[rank] (r2) at (2.72,0) {$r_2$};
  \node[rank] (dots) at (4.08,0) {$\cdots$};
  \node[rank] (rq) at (5.44,0) {$r_q$};

  \node[count] (drop) at (2.72,-0.84) {$w_p=r_{p-1}-r_p$};
  \node[count, minimum width=3.28cm] (exact) at (2.72,-1.55) {size-$p$ blocks $=w_p-w_{p+1}$};

  \draw[arrow] (r0) -- (r1);
  \draw[arrow] (r1) -- (r2);
  \draw[arrow] (r2) -- (dots);
  \draw[arrow] (dots) -- (rq);
  \draw[arrow] (drop) -- (exact);
\end{tikzpicture}

  \caption{For $r_p=\operatorname{rank}((A-\lambda I)^p)$, drops count surviving chains; $w_p-w_{p+1}$ counts size-$p$ blocks.}
  \label{fig:weyr-rank-drops-jordan-blocks}
\end{figure}

\begin{method} \label{meth:jordan-form-from-rank-drops} \hlt{Computing Jordan Blocks by Rank Drops}\\
For each eigenvalue $\lambda$ of $A\in\operatorname{Mat}(n,n,\mathbb{C})$:
\begin{enumerate}[label=(\roman*)]
\setlength{\itemsep}{0pt}
\item Compute $r_p=\operatorname{rank}\bigl((A-\lambda I_n)^p\bigr)$ for $p=0,1,2,\ldots$ until the differences stabilize at $0$.
\item Set $w_p=r_{p-1}-r_p$. Then $w_p$ is the number of $\lambda$-blocks of size at least $p$.
\item The number of blocks $J_p(\lambda)$ is $w_p-w_{p+1}$.
\item Assemble the Jordan form by taking the direct sum over all eigenvalues.
\end{enumerate}
\end{method}

\begin{theorem} \label{thm:ab-ba-nonzero-eigenvalues} \hlt{Nonzero Eigenvalues of $AB$ and $BA$}\\
Let $A\in\operatorname{Mat}(m,n,\mathbb{C})$, $B\in\operatorname{Mat}(n,m,\mathbb{C})$, and suppose $n\geq m$. Then $p_{BA}(z)=z^{n-m}p_{AB}(z)$.\\
Consequently, for $\lambda\neq0$, $m_{AB}(\lambda)=m_{BA}(\lambda)$. If $m_{AB}(0)=k$, then $m_{BA}(0)=k+n-m$.
\begin{proof}
For $z\neq0$, Sylvester's determinant identity gives
\begin{align}
\det(zI_n-BA) =z^n\det(I_n-z^{-1}BA) =z^n\det(I_m-z^{-1}AB) =z^{n-m}\det(zI_m-AB).
\nonumber
\end{align}
Both sides are polynomials in $z$, so the identity holds for all $z$. The claimed multiplicities follow by factoring $p_{AB}$ and comparing with $p_{BA}=z^{n-m}p_{AB}$.\\
\hlt{Structural Commentary}: Sylvester's determinant identity moves the rectangular product across the determinant.\\
\hlt{Key Insight}: Passing from $AB$ to $BA$ only adds the necessary extra zero roots when the sizes differ.
\end{proof}
\end{theorem}

\begin{theorem} \label{thm:ab-ba-same-nonzero-jordan-blocks} \hlt{Nonzero Jordan Blocks of $AB$ and $BA$}\\
Let $A,B\in\operatorname{Mat}(n,n,\mathbb{C})$. For every $\lambda\neq0$, the Jordan blocks of $AB$ and $BA$ at $\lambda$ coincide.
\begin{proof}
Set $X=\begin{bmatrix}AB&A\\0&0\end{bmatrix}$ and $Y=\begin{bmatrix}0&A\\0&BA\end{bmatrix}$. Then $\begin{bmatrix}I_n&0\\B&I_n\end{bmatrix}
X
\begin{bmatrix}I_n&0\\-B&I_n\end{bmatrix}
=Y$, so $X$ and $Y$ are similar.\\
For $\lambda\neq0$ and $p\geq1$, the shifted powers have block triangular forms
\begin{align}
(X-\lambda I_{2n})^p=
\begin{bmatrix}
(AB-\lambda I_n)^p&*\\
0&(-\lambda)^pI_n
\end{bmatrix},
\qquad
(Y-\lambda I_{2n})^p=
\begin{bmatrix}
(-\lambda)^pI_n&*\\
0&(BA-\lambda I_n)^p
\end{bmatrix} \nonumber
\end{align}
The diagonal block $(-\lambda)^pI_n$ is invertible.\\
Block Gaussian elimination by invertible row and column operations removes the starred block, so
\begin{align}
\operatorname{rank}\bigl((X-\lambda I_{2n})^p\bigr)
&=n+\operatorname{rank}\bigl((AB-\lambda I_n)^p\bigr), \nonumber\\
\operatorname{rank}\bigl((Y-\lambda I_{2n})^p\bigr)
&=n+\operatorname{rank}\bigl((BA-\lambda I_n)^p\bigr).
\nonumber
\end{align}
Since $X$ and $Y$ are similar, these left sides agree for all $p$.\\
Hence the shifted rank drops of $AB$ and $BA$ at $\lambda$ agree, so Theorem \ref{thm:jordan-block-sizes-from-rank-drops} gives the same Jordan blocks.\\
\hlt{Structural Commentary}: Similar block embeddings make the nonzero shifted ranks match.\\
\hlt{Key Insight}: At $\lambda\neq0$, the extra block is invertible and splits off.
\end{proof}
\end{theorem}

\begin{figure}[H]
  \centering
  \begin{tikzpicture}[
  font=\small,
  box/.style={draw=blue!60!black, fill=blue!6, rounded corners=2pt, align=center, inner sep=6pt, minimum width=2.9cm},
  aux/.style={draw=gray!65, fill=gray!6, rounded corners=2pt, align=center, inner sep=6pt, minimum width=3.2cm},
  arr/.style={-{Stealth[length=2mm]}, thick, gray!75}
]
  \node[aux] (x) at (0,0) {$X-\lambda I_{2n}$};
  \node[aux] (y) at (4.8,0) {$Y-\lambda I_{2n}$};
  \node[box] (ab) at (0,-1.55) {$AB-\lambda I_n$};
  \node[box] (ba) at (4.8,-1.55) {$BA-\lambda I_n$};

  \draw[arr, <->] (x) -- node[above] {$\sim$} (y);
  \draw[arr] (x) -- (ab);
  \draw[arr] (y) -- (ba);
\end{tikzpicture}

  \caption{For $\lambda\neq0$, the auxiliary similarity reduces the comparison to shifted ranks of $AB$ and $BA$.}
  \label{fig:ab-ba-nonzero-jordan-blocks}
\end{figure}

\begin{corollary} \label{cor:jordan-form-unique-up-to-block-permutation} \hlt{Uniqueness of Jordan Form}\\
If $J$ and $J'$ are Jordan matrices similar to $A\in\operatorname{Mat}(n,n,\mathbb{C})$, they differ only by a permutation of Jordan blocks.
\begin{proof}
Similarity preserves spectra and the ranks of $(X-\lambda I_n)^p$. Thus $J$ and $J'$ have the same Weyr characteristic for each eigenvalue. Theorem \ref{thm:jordan-block-sizes-from-rank-drops} then gives the same block multiplicities, so only block order may differ.\\
\hlt{Structural Commentary}: Weyr data is similarity-invariant and fixes block multiplicities.\\
\hlt{Key Insight}: Jordan form is unique modulo block order.
\end{proof}
\end{corollary}

\begin{corollary} \label{cor:ab-ba-similar-rank-power-criterion} \hlt{Similarity Criterion for $AB$ and $BA$}\\
Let $A,B\in\operatorname{Mat}(n,n,\mathbb{C})$. $AB$ is similar to $BA$ if and only if $\operatorname{rank}\bigl((AB)^p\bigr)=\operatorname{rank}\bigl((BA)^p\bigr)$, $p=1,\ldots,n$.
\begin{proof}
\begin{enumerate}[label=\roman*.]
\setlength{\itemsep}{0pt}
\item[$(\Rightarrow)$]
If $AB\sim BA$, then $(AB)^p\sim(BA)^p$ for all $p$, so their ranks are equal.
\item[$(\Leftarrow)$]
Assume the rank equalities hold for $p=1,\ldots,n$. Theorem \ref{thm:ab-ba-same-nonzero-jordan-blocks} gives the same nonzero Jordan blocks. Together with the equality at $p=0$, the rank equalities give the same Weyr characteristic at $\lambda=0$ by Theorem \ref{thm:jordan-block-sizes-from-rank-drops}. Hence the nilpotent blocks also agree, so the Jordan forms coincide up to block order.
\end{enumerate}
\hlt{Structural Commentary}: Nonzero blocks agree automatically; only the nilpotent blocks must be tested.\\
\hlt{Key Insight}: Powers of $AB$ and $BA$ measure exactly the zero-Jordan-block data.
\end{proof}
\end{corollary}

\begin{method} \label{meth:checking-ab-ba-similarity} \hlt{Checking Whether $AB$ and $BA$ Are Similar}\\
For $A,B\in\operatorname{Mat}(n,n,\mathbb{C})$:
\begin{enumerate}[label=(\roman*)]
\setlength{\itemsep}{0pt}
\item The nonzero Jordan blocks of $AB$ and $BA$ already match.
\item Compute $\operatorname{rank}\bigl((AB)^p\bigr)$ and $\operatorname{rank}\bigl((BA)^p\bigr)$ for $p=1,\ldots,n$.
\item If all ranks agree, then $AB\sim BA$; if some rank differs, then only the nilpotent Jordan structure differs.
\end{enumerate}
\end{method}

\begin{definition} \label{def:nonderogatory-matrix} \hlt{Nonderogatory Matrix}\\
A matrix $A\in\operatorname{Mat}(n,n,\mathbb{C})$ is \hlt{nonderogatory} if its Jordan form has exactly one Jordan block for each distinct eigenvalue. Equivalently, $\dim E_{\lambda}(A)=1$ for every $\lambda\in\operatorname{spec}(A)$.
\end{definition}

\begin{lemma} \label{lem:single-jordan-block-commutant-polynomials} \hlt{Commutant of One Jordan Block}\\
If $J=J_m(\lambda)$ and $XJ=JX$, then $X=p(J)$ for some polynomial $p$ with $\deg p<m$.
\begin{proof}
It suffices to treat $N=J_m(0)$, since $J_m(\lambda)=\lambda I_m+N$.\\
The vector $\mathbf{e}_m$ is cyclic for $N$: $(\mathbf{e}_m,N\mathbf{e}_m,\ldots,N^{m-1}\mathbf{e}_m)$ is the standard basis in reverse order.\\
Choose the unique polynomial $p$ with $\deg p<m$ such that $p(N)\mathbf{e}_m=X\mathbf{e}_m$.\\
For every $r=0,\ldots,m-1$, $p(N)N^r\mathbf{e}_m=N^rp(N)\mathbf{e}_m=N^rX\mathbf{e}_m=XN^r\mathbf{e}_m$.\\
Thus $p(N)$ and $X$ agree on a basis, so $X=p(N)=q(J)$ for $q(z)=p(z-\lambda)$.\\
\hlt{Structural Commentary}: A single Jordan block is generated by one cyclic vector.\\
\hlt{Key Insight}: The cyclic vector determines the commutant.
\end{proof}
\end{lemma}

\begin{theorem} \label{thm:nonderogatory-commutant-polynomials} \hlt{Nonderogatory Commutants Are Polynomials}\\
If $A\in\operatorname{Mat}(n,n,\mathbb{C})$ is nonderogatory, then $\operatorname{Comm}(A)=\{p(A):p\in\mathbb{C}[z]\}$.
\begin{proof}
The inclusion $\{p(A):p\in\mathbb{C}[z]\}\subseteq\operatorname{Comm}(A)$ follows from Proposition \ref{prop:operator-polynomial-products}.\\
For reverse inclusion, first replace $A$ by its Jordan form $J=J_{m_1}(\lambda_1)\oplus\cdots\oplus J_{m_d}(\lambda_d)$, where $\lambda_i$ are distinct.\\
If $XJ=JX$ and $X=[X_{ij}]$ is partitioned conformally with $J$, the block equation is $X_{ij}J_{m_j}(\lambda_j)=J_{m_i}(\lambda_i)X_{ij}$. Corollary \ref{cor:commuting-with-disjoint-blocks-has-zero-off-diagonal-blocks} gives $X_{ij}=0$ for $i\neq j$. Each diagonal block $X_{ii}$ commutes with $J_{m_i}(\lambda_i)$, so Lemma \ref{lem:single-jordan-block-commutant-polynomials} gives $X_{ii}=p_i(J_{m_i}(\lambda_i))$ for some polynomial $p_i$ with $\deg p_i<m_i$. The ideals generated by $(z-\lambda_i)^{m_i}$ are pairwise coprime, so Hermite interpolation gives a polynomial $p$ with $p\equiv p_i\pmod{(z-\lambda_i)^{m_i}}$ for every $i$.\\
On the $i$th block, $p-p_i=(z-\lambda_i)^{m_i}q_i$ for some $q_i$, and $(J_{m_i}(\lambda_i)-\lambda_iI)^{m_i}=0$. Hence $p(J_{m_i}(\lambda_i))=p_i(J_{m_i}(\lambda_i))=X_{ii}$ for every $i$, so $X=p(J)$. Conjugating back to $A$ proves the claim.\\
\hlt{Structural Commentary}: Disjoint eigenvalue blocks force block diagonality; one Jordan block per eigenvalue makes each diagonal commutant polynomial.\\
\hlt{Key Insight}: Hermite interpolation glues the separate block polynomials into one polynomial in $A$.
\end{proof}
\end{theorem}

\begin{figure}[H]
  \centering
  \begin{tikzpicture}[
  font=\small,
  block/.style={draw=blue!60!black, fill=blue!6, rounded corners=2pt, align=center, inner sep=5pt, minimum width=2.00cm},
  poly/.style={draw=red!60!black, fill=red!6, rounded corners=2pt, align=center, inner sep=5pt, minimum width=2.33cm},
  arr/.style={-{Stealth[length=2mm]}, thick, gray!75}
]
  \node[block] (j1) at (0,0) {$J_{m_1}(\lambda_1)$};
  \node[block] (j2) at (2.55,0) {$J_{m_2}(\lambda_2)$};
  \node[block] (jd) at (5.10,0) {$\cdots$};
  \node[block] (jk) at (7.65,0) {$J_{m_d}(\lambda_d)$};

  \node[poly] (p1) at (0,-1.30) {$p_1(J_{m_1}(\lambda_1))$};
  \node[poly] (p2) at (2.55,-1.30) {$p_2(J_{m_2}(\lambda_2))$};
  \node[poly] (pd) at (5.10,-1.30) {$\cdots$};
  \node[poly] (pk) at (7.65,-1.30) {$p_d(J_{m_d}(\lambda_d))$};

  \node[poly, minimum width=3.52cm] (p) at (3.825,-2.52) {$p\equiv p_i\pmod{(z-\lambda_i)^{m_i}}$};

  \draw[arr] (j1) -- (p1);
  \draw[arr] (j2) -- (p2);
  \draw[arr] (jk) -- (pk);
  \draw[arr] (p1) -- (p);
  \draw[arr] (p2) -- (p);
  \draw[arr] (pk) -- (p);
\end{tikzpicture}

  \caption{Blockwise commutant polynomials glue by interpolation.}
  \label{fig:nonderogatory-commutant-polynomial}
\end{figure}

\begin{theorem} \label{thm:conjugate-similarity-jordan-block-criterion} \hlt{Conjugate Similarity Criterion}\\
Let $A\in\operatorname{Mat}(n,n,\mathbb{C})$. Then $A$ is similar to $\overline{A}$ if and only if each $J_k(\lambda)$ occurs in the Jordan form of $A$ with the same multiplicity as $J_k(\overline{\lambda})$.
\begin{proof}
Let $J$ be a Jordan form of $A$. Then $\overline{J}$ is a Jordan form of $\overline{A}$, obtained by replacing each $J_k(\lambda)$ by $J_k(\overline{\lambda})$.
\begin{enumerate}[label=\roman*.]
\setlength{\itemsep}{0pt}
\item[$(\Rightarrow)$] If $A$ is similar to $\overline{A}$, then $J$ and $\overline{J}$ are similar.\\
By Corollary \ref{cor:jordan-form-unique-up-to-block-permutation}, their Jordan blocks have same multiplicities.
\item[$(\Leftarrow)$] If the multiplicities agree, then $J$ and $\overline{J}$ differ only by a permutation of blocks. Hence $A$ is similar to $\overline{A}$.
\end{enumerate}
\hlt{Structural Commentary}: Conjugation changes eigenvalue labels but not block sizes.\\
\hlt{Key Insight}: Non-real Jordan blocks must occur in conjugate pairs.
\end{proof}
\end{theorem}

\begin{corollary} \label{cor:hermitian-products-and-conjugate-products} \hlt{Hermitian Products and Conjugate Products}\\
Let $A,H,K\in\operatorname{Mat}(n,n,\mathbb{C})$, and suppose $H$ and $K$ are Hermitian.\\
Then $HK$ is similar to $KH$, and $A\overline{A}$ is similar to $\overline{A}A$.
\begin{proof}
For each $p=1,\ldots,n$, adjoints preserves rank, $((HK)^p)^*=(KH)^p$. Hence $\operatorname{rank}\bigl((HK)^p\bigr)=\operatorname{rank}\bigl((KH)^p\bigr)$.\\
Similarly, entrywise conjugation preserves rank and gives $\overline{(A\overline{A})^p}=(\overline{A}A)^p$, so $\operatorname{rank}\bigl((A\overline{A})^p\bigr)=\operatorname{rank}\bigl((\overline{A}A)^p\bigr)$.\\
Apply Corollary \ref{cor:ab-ba-similar-rank-power-criterion} in both cases.\\
\hlt{Structural Commentary}: Adjoint and conjugation preserve rank and reverse product order.\\
\hlt{Key Insight}: The rank-power test turns order reversal into similarity.
\end{proof}
\end{corollary}

\begin{definition} \label{def:reversal-matrix} \hlt{Reversal Matrix}\\
For $k\geq1$, \hlt{reversal matrix} $R_k\in\operatorname{Mat}(k,k,\mathbb{C})$ defined $(R_k)_{ij}=
\begin{cases}
1,& i+j=k+1,\\
0,& i+j\neq k+1
\end{cases}$, thus $R_k^t=R_k$, $R_k^2=I_k$.
\end{definition}

\begin{lemma} \label{lem:jordan-block-transpose-reversal-similarity} \hlt{Reversal Conjugates a Jordan Block to Its Transpose}\\
For every $k\geq1$ and $\lambda\in\mathbb{C}$, $J_k(\lambda)^t=R_kJ_k(\lambda)R_k$. Moreover, $R_kJ_k(\lambda)$ and $J_k(\lambda)R_k$ are symmetric.
\begin{proof}
Let $N=J_k(0)$. Since $R_k$ reverses standard basis order, $R_kNR_k=N^t$.\\
Adding scalar gives $R_kJ_k(\lambda)R_k=J_k(\lambda)^t$. As $R_k=R_k^t=R_k^{-1}$, this identity gives $J_k(\lambda)^tR_k=R_kJ_k(\lambda)$ and $R_kJ_k(\lambda)^t=J_k(\lambda)R_k$, so $R_kJ_k(\lambda)$ and $J_k(\lambda)R_k$ are symmetric.\\
\hlt{Structural Commentary}: Reversal changes the superdiagonal shift into the subdiagonal shift.\\
\hlt{Key Insight}: Transposition reads the same Jordan chain backwards.
\end{proof}
\end{lemma}

\begin{theorem} \label{thm:matrix-similar-to-transpose-and-symmetric-factorizations} \hlt{Transpose Similarity and Symmetric Factorisations}\\
Let $A\in\operatorname{Mat}(n,n,\mathbb{C})$. Then:
\begin{enumerate}[label=(\roman*)]
\setlength{\itemsep}{0pt}
\item $A$ is similar to $A^t$ via a symmetric similarity matrix.
\item $A=BC=DE$, where $B,C,D,E$ are symmetric and $B,E$ are invertible.
\end{enumerate}
\begin{proof}
Choose a Jordan form $A=SJS^{-1}$, and let $R$ be the direct sum of the reversal matrices matching the Jordan blocks of $J$. By Lemma \ref{lem:jordan-block-transpose-reversal-similarity}, $J^t=RJR$, $R^t=R$, $R^2=I_n$, and both $RJ$ and $JR$ are symmetric.
\begin{enumerate}[label=(\roman*)]
\setlength{\itemsep}{0pt}
\item
Set $P=(S^{-1})^tRS^{-1}$. Then $P^t=P$, $P^{-1}=SRS^t$, and $PAP^{-1}=(S^{-1})^tRJR S^t=(S^{-1})^tJ^tS^t=A^t$.
\item
Using $R^2=I_n$, $A=(SRS^t)\bigl((S^{-1})^tRJ S^{-1}\bigr)$ and $A=(SJR S^t)\bigl((S^{-1})^tR S^{-1}\bigr)$.\\
The products equal $SJS^{-1}$ because $S^t(S^{-1})^t=I_n$ and $R^2=I_n$. All four factors are symmetric because $R,RJ,JR$ are symmetric; $SRS^t$ and $(S^{-1})^tRS^{-1}$ are invertible.
\end{enumerate}
\hlt{Structural Commentary}: Reduce to Jordan blocks, reverse chains, then conjugate back.\\
\hlt{Key Insight}: Symmetric involutions produce both transpose similarity and symmetric factors.
\end{proof}
\end{theorem}

\begin{figure}[H]
  \centering
  \begin{tikzpicture}[
  font=\small,
  vec/.style={draw=blue!65!black, fill=blue!7, rounded corners=2pt, minimum height=0.50cm, minimum width=0.72cm, align=center},
  rev/.style={draw=red!60!black, fill=red!7, rounded corners=2pt, minimum height=0.50cm, minimum width=0.72cm, align=center},
  arr/.style={-{Stealth[length=2mm]}, thick, blue!70!black},
  rarr/.style={-{Stealth[length=2mm]}, thick, red!70!black},
  flip/.style={gray!65, dashed}
]
  \node[vec] (e1) at (0,0) {$\mathbf{e}_1$};
  \node[vec] (e2) at (1.06,0) {$\mathbf{e}_2$};
  \node[vec] (ed) at (2.12,0) {$\cdots$};
  \node[vec] (ek) at (3.18,0) {$\mathbf{e}_k$};

  \draw[arr] (ek) -- (ed);
  \draw[arr] (ed) -- (e2);
  \draw[arr] (e2) -- (e1);
  \draw[arr] (e1) -- ++(-0.74,0) node[left] {$0$};

  \node[rev] (r1) at (0,-1.00) {$\mathbf{e}_k$};
  \node[rev] (r2) at (1.06,-1.00) {$\mathbf{e}_{k-1}$};
  \node[rev] (rd) at (2.12,-1.00) {$\cdots$};
  \node[rev] (rk) at (3.18,-1.00) {$\mathbf{e}_1$};

  \draw[rarr] (r1) -- (r2);
  \draw[rarr] (r2) -- (rd);
  \draw[rarr] (rd) -- (rk);
  \draw[rarr] (rk) -- ++(0.74,0) node[right] {$0$};

  \draw[flip] (e1) -- (r1);
  \draw[flip] (ek) -- (rk);
\end{tikzpicture}

  \caption{Reversal gives $R_kJ_k(0)R_k=J_k(0)^t$: the same chain read backwards.}
  \label{fig:jordan-transpose-reversal}
\end{figure}

\begin{theorem} \label{thm:conjugate-similarity-hermitian-product-equivalences} \hlt{Conjugate Similarity and Hermitian Products}\\
Let $A\in\operatorname{Mat}(n,n,\mathbb{C})$. The following are equivalent:
\begin{enumerate}[label=(\roman*)]
\setlength{\itemsep}{0pt}
\item $A$ is similar to $\overline{A}$.
\item $A$ is similar to $A^*$.
\item $A$ is similar to $A^*$ via a Hermitian similarity matrix.
\item $A=HK=LM$, where $H,K,L,M$ are Hermitian, and $H,M$ are invertible.
\item $A=HK$ for some Hermitian matrices $H,K$.
\end{enumerate}
\begin{proof}
\begin{enumerate}[label=\roman*.]
\setlength{\itemsep}{0pt}
\item[$(i)$]$\Rightarrow(ii)$ Theorem \ref{thm:matrix-similar-to-transpose-and-symmetric-factorizations} applied to $\overline{A}$ gives $\overline{A}\sim(\overline{A})^t=A^*$. Thus $A\sim A^*$.
\item[$(ii)$]$\Rightarrow(iii)$ Choose invertible $S$ such that $A=SA^*S^{-1}$, so $AS=SA^*$. For $\theta\in\mathbb{R}$, put $S_{\theta}=e^{i\theta}S$.\\
Then $AS_{\theta}=S_{\theta}A^*$ and, by taking adjoints, $AS_{\theta}^*=S_{\theta}^*A^*$. Hence $A(S_{\theta}+S_{\theta}^*)=(S_{\theta}+S_{\theta}^*)A^*$.\\
The matrix $H_{\theta}=S_{\theta}+S_{\theta}^*$ is Hermitian. Also $S_{\theta}^{-1}H_{\theta}=I_n+e^{-2i\theta}S^{-1}S^*$. Thus $H_{\theta}$ is singular exactly when $-e^{2i\theta}$ is an eigenvalue of $S^{-1}S^*$. There are only finitely many such values, so choose $\theta$ for which $H_{\theta}$ is invertible. Then $A=H_{\theta}A^*H_{\theta}^{-1}$.
\item[$(iii)$]$\Rightarrow(iv)$ Let $R$ be invertible Hermitian and satisfy $A=RA^*R^{-1}$. Then $AR=RA^*$ and $R^{-1}A=A^*R^{-1}$.\\
Thus $A^*R^{-1}$ and $RA^*$ are Hermitian, and $A=R(A^*R^{-1})=(RA^*)R^{-1}$.\\
This gives the two asserted Hermitian factorizations, with the first factor in the first product and the second factor in the second product invertible.
\item[$(iv)$]$\Rightarrow(v)$ The first factorization in (iv) already has the form $A=HK$ with $H$ and $K$ Hermitian.
\item[$(v)$]$\Rightarrow(i)$ If $A=HK$ with $H,K$ Hermitian, then Corollary \ref{cor:hermitian-products-and-conjugate-products} gives $HK\sim KH$.\\
By Theorem \ref{thm:matrix-similar-to-transpose-and-symmetric-factorizations}, $KH\sim (KH)^t=H^tK^t=\overline{H}\,\overline{K}=\overline{A}$. Hence $A\sim\overline{A}$.
\end{enumerate}
\hlt{Structural Commentary}: The equivalences move between conjugation, adjoints, and Hermitian products by combining transpose similarity with a Hermitianized similarity matrix.\\
\hlt{Key Insight}: A phase-rotated average $e^{i\theta}S+e^{-i\theta}S^*$ turns an arbitrary similarity to $A^*$ into a Hermitian one.
\end{proof}
\end{theorem}

\begin{figure}[H]
  \centering
  \begin{tikzpicture}[
  font=\small,
  node/.style={draw=blue!60!black, fill=blue!6, rounded corners=2pt, align=center, inner sep=6pt, minimum width=2.9cm},
  mid/.style={draw=red!60!black, fill=red!6, rounded corners=2pt, align=center, inner sep=6pt, minimum width=3.3cm},
  arr/.style={-{Stealth[length=2mm]}, thick, gray!75},
  lbl/.style={fill=white, text=gray!70!black, inner xsep=2pt, inner ysep=1pt}
]
  \node[node] (conj) at (0,0) {$A\sim\overline{A}$};
  \node[node] (adj) at (5.0,0) {$A\sim A^*$};
  \node[mid] (herm) at (10.4,0) {$A=HA^*H^{-1}$\\$H=H^*$};
  \node[node] (prod) at (5.0,-2.2) {$A=HK$\\$H=H^*,\ K=K^*$};

  \draw[arr] (conj.east) -- node[lbl, above, yshift=2pt] {transpose} (adj.west);
  \draw[arr] (adj.east) -- node[lbl, above, yshift=2pt] {phase average} (herm.west);
  \draw[arr] (herm.south west) -- node[lbl, right, xshift=2pt] {factor} (prod.east);
  \draw[arr] (prod.west) -- node[lbl, left, xshift=-2pt] {$HK\sim KH$} (conj.south east);
\end{tikzpicture}

  \caption{Conjugate similarity, adjoint similarity, and Hermitian product factorisations are equivalent.}
  \label{fig:conjugate-similarity-hermitian-products}
\end{figure}

\begin{definition} \label{def:complexification-real-vector-space} \hlt{Complexification of a Real Vector Space}\\
Let $\mathcal{V}$ be a real vector space. The \hlt{complexification} of $\mathcal{V}$ is $\mathcal{V}_{\mathbb{C}}\defeq \mathcal{V}\times\mathcal{V}$, where $(\mathbf{u},\mathbf{v})$ is written as $\mathbf{u}+i\mathbf{v}$.\\
Addition and complex scalar multiplication are defined by
\begin{align}
(\mathbf{u}_1+i\mathbf{v}_1)+(\mathbf{u}_2+i\mathbf{v}_2)
&=(\mathbf{u}_1+\mathbf{u}_2)+i(\mathbf{v}_1+\mathbf{v}_2), \nonumber\\
(a+bi)(\mathbf{u}+i\mathbf{v})
&=(a\mathbf{u}-b\mathbf{v})+i(b\mathbf{u}+a\mathbf{v}),
\qquad a,b\in\mathbb{R}.
\nonumber
\end{align}
We identify $\mathcal{V}$ with the real subspace $\{\mathbf{u}+i\mathbf{0}:\mathbf{u}\in\mathcal{V}\}$.
\end{definition}

\begin{proposition} \label{prop:complexification-vector-space-basis} \hlt{Complexification Preserves Bases}\\
Let $\mathcal{V}$ be real. Then $\mathcal{V}_{\mathbb{C}}$ is complex, and every real basis $(\mathbf{v}_1,\ldots,\mathbf{v}_n)$ of $\mathcal{V}$ is a complex basis of $\mathcal{V}_{\mathbb{C}}$.\\
Thus $\dim_{\mathbb{C}}(\mathcal{V}_{\mathbb{C}})=\dim_{\mathbb{R}}(\mathcal{V})$.
\begin{proof}
The vector-space axioms follow componentwise from those of $\mathcal{V}$.\\
For $\mathbf{u},\mathbf{v}\in\mathcal{V}$, write $\mathbf{u}=\sum\limits_{j=1}^{n}a_j\mathbf{v}_j$ and $\mathbf{v}=\sum\limits_{j=1}^{n}b_j\mathbf{v}_j$, $a_j,b_j\in\mathbb{R}$. Then $\mathbf{u}+i\mathbf{v}=\sum\limits_{j=1}^{n}(a_j+ib_j)\mathbf{v}_j$, so the list spans.\\
If $\sum\limits_{j=1}^{n}\lambda_j\mathbf{v}_j=\mathbf{0}$ and $\lambda_j=a_j+ib_j$, then $\sum\limits_{j=1}^{n}a_j\mathbf{v}_j=\mathbf{0}$ and $\sum\limits_{j=1}^{n}b_j\mathbf{v}_j=\mathbf{0}$.\\
Real linear independence gives $a_j=b_j=0$ for all $j$, hence $\lambda_j=0$ for all $j$.\\
\hlt{Structural Commentary}: Complexification keeps the basis length unchanged.\\
\hlt{Key Insight}: Separate real and imaginary parts.
\end{proof}
\end{proposition}

\begin{proposition} \label{prop:complexification-real-dimension-and-embedding} \hlt{Complexification as a Real Space}\\
If $\mathcal{V}$ has real basis $(\mathbf{v}_1,\ldots,\mathbf{v}_n)$, then $\mathcal{V}_{\mathbb{C}}$, as a real vector space, has basis $(\mathbf{v}_1,\ldots,\mathbf{v}_n,i\mathbf{v}_1,\ldots,i\mathbf{v}_n)$.\\
Hence $\dim_{\mathbb{R}}(\mathcal{V}_{\mathbb{C}})=2\dim_{\mathbb{R}}(\mathcal{V})$, and $\operatorname{cpx}:\mathcal{V}\to\mathcal{V}_{\mathbb{C}}$, $\operatorname{cpx}(\mathbf{v})=\mathbf{v}+i\mathbf{0}$, is real-linear and injective.
\begin{proof}
For $\mathbf{u}=\sum a_j\mathbf{v}_j$ and $\mathbf{w}=\sum b_j\mathbf{v}_j$, $\mathbf{u}+i\mathbf{w}=\sum a_j\mathbf{v}_j+\sum b_j(i\mathbf{v}_j)$, so the displayed list spans $\mathcal{V}_{\mathbb{C}}$ over $\mathbb{R}$.\\
If $\sum a_j\mathbf{v}_j+\sum b_j(i\mathbf{v}_j)=\mathbf{0}$, then $\operatorname{Re}$ and $\operatorname{Im}$ give $\sum a_j\mathbf{v}_j=\sum b_j\mathbf{v}_j=\mathbf{0}$; all $a_j,b_j$ vanish. Thus this is real basis.\\
For $r\in\mathbb{R}$, $\operatorname{cpx}(r\mathbf{v}+\mathbf{w})=r\operatorname{cpx}(\mathbf{v})+\operatorname{cpx}(\mathbf{w})$, and $\ker(\operatorname{cpx})=\{\mathbf{0}\}$, so $\operatorname{cpx}$ is real-linear and injective.\\
\hlt{Structural Commentary}: Complexification is two real copies of $\mathcal{V}$.\\
\hlt{Key Insight}: The real and imaginary parts are independent real copies of the same basis.
\end{proof}
\end{proposition}

\begin{definition} \label{def:complexification-real-operator} \hlt{Complexification of a Real Operator}\\
Let $\mathcal{V}$ be a real vector space and let $T\in\LL(\mathcal{V})$.\\
The \hlt{complexification} of $T$ is the operator $T_{\mathbb{C}}\in\LL(\mathcal{V}_{\mathbb{C}})$ defined by $T_{\mathbb{C}}(\mathbf{u}+i\mathbf{v})\defeq T\mathbf{u}+iT\mathbf{v}$.
\end{definition}

\begin{proposition} \label{prop:complexification-operator-basic-properties} \hlt{Complexification Preserves Operator Data}\\
Let $\mathcal{V}$ be a real vector space and let $T\in\LL(\mathcal{V})$.
\begin{enumerate}[label=(\roman*)]
\setlength{\itemsep}{0pt}
\item $T_{\mathbb{C}}\in\LL(\mathcal{V}_{\mathbb{C}})$.
\item $(T_{\mathbb{C}})^k(\mathbf{u}+i\mathbf{v})=T^k\mathbf{u}+iT^k\mathbf{v}$ for every integer $k\geq0$.
\item If $\beta$ is a real basis of $\mathcal{V}$, then $[T_{\mathbb{C}}]_{\beta\leftarrow\beta}=[T]_{\beta\leftarrow\beta}$, now regarded as a complex matrix.
\end{enumerate}
\begin{proof}
\begin{enumerate}[label=(\roman*)]
\setlength{\itemsep}{0pt}
\item Additivity is immediate. For $\alpha=a+bi\in\mathbb{C}$,
\begin{align}
T_{\mathbb{C}}\bigl(\alpha(\mathbf{u}+i\mathbf{v})\bigr) =T(a\mathbf{u}-b\mathbf{v})+iT(b\mathbf{u}+a\mathbf{v}) =(aT\mathbf{u}-bT\mathbf{v})+i(bT\mathbf{u}+aT\mathbf{v}) =\alpha T_{\mathbb{C}}(\mathbf{u}+i\mathbf{v}),
\nonumber
\end{align}
so $T_{\mathbb{C}}$ is complex-linear.
\item For $k=0$, the claim is the identity. If it holds for $k$, then
\begin{align}
(T_{\mathbb{C}})^{k+1}(\mathbf{u}+i\mathbf{v}) =T_{\mathbb{C}}(T^k\mathbf{u}+iT^k\mathbf{v}) =T^{k+1}\mathbf{u}+iT^{k+1}\mathbf{v}.
\nonumber
\end{align}
Thus the formula holds for all $k\geq0$.
\item If $T\mathbf{v}_j=\sum\limits_i a_{ij}\mathbf{v}_i$ with $a_{ij}\in\mathbb{R}$, then $T_{\mathbb{C}}\mathbf{v}_j=\sum\limits_i a_{ij}\mathbf{v}_i$. Hence the coordinate matrix is unchanged.
\end{enumerate}
\hlt{Structural Commentary}: Complexification keeps the coordinate matrix unchanged.\\
\hlt{Key Insight}: Same real entries; larger scalar field.
\end{proof}
\end{proposition}

\begin{figure}[H]
  \centering
  \begin{tikzpicture}[
  >=Stealth,
  every node/.style={font=\small},
  box/.style={draw=gray!55, rounded corners=3pt, inner xsep=7pt, inner ysep=4.5pt, align=center},
  arrow/.style={->, line width=0.85pt, blue!70!black}
]
  \node[box] (real) at (0,0) {$\mathbf{u},\mathbf{v}\in\mathcal{V}$};
  \node[box] (complex) at (2.65,0) {$\mathbf{u}+i\mathbf{v}\in\mathcal{V}_{\mathbb{C}}$};
  \node[box] (image) at (5.80,0) {$T\mathbf{u}+iT\mathbf{v}$};

  \draw[arrow] (real) -- (complex);
  \draw[arrow] (complex) -- (image);
\end{tikzpicture}

  \caption{Complexification packages two real vectors: $T_{\mathbb{C}}(\mathbf{u}+i\mathbf{v})=T\mathbf{u}+iT\mathbf{v}$.}
  \label{fig:complexification-real-operator}
\end{figure}

\begin{method} \label{meth:complexify-real-matrix} \hlt{Complexifying a Real Matrix}\\
To use complex eigenvalues for a real matrix $A\in\operatorname{Mat}(n,n,\mathbb{R})$:
\begin{enumerate}[label=(\roman*)]
\setlength{\itemsep}{0pt}
\item Regard the same real matrix $A$ as a complex matrix on $\mathbb{C}^n$.
\item Find complex eigenpairs or invariant data for this complexified operator.
\item If $A(\mathbf{u}+i\mathbf{v})=(a+bi)(\mathbf{u}+i\mathbf{v})$, separate parts: $A\mathbf{u}=a\mathbf{u}-b\mathbf{v}$, $A\mathbf{v}=b\mathbf{u}+a\mathbf{v}$
\item Conclude that $\operatorname{span}(\mathbf{u},\mathbf{v})\subseteq\mathbb{R}^n$ is an invariant line or plane.
\end{enumerate}
\end{method}

\begin{theorem} \label{thm:real-operator-invariant-line-or-plane} \hlt{Invariant Line or Plane over $\mathbb{R}$}\\
Every operator on a nonzero finite-dimensional real vector space has an invariant subspace of dimension $1$ or $2$.
\begin{proof}
Let $\mathcal{V}$ be a nonzero finite-dimensional real vector space and let $T\in\LL(\mathcal{V})$.\\
By Theorem \ref{thm:complex-operator-has-eigenvalue}, $T_{\mathbb{C}}$ has an eigenvalue $\lambda=a+bi$ with eigenvector $\mathbf{u}+i\mathbf{v}\neq\mathbf{0}$.\\
The equation $T_{\mathbb{C}}(\mathbf{u}+i\mathbf{v})=(a+bi)(\mathbf{u}+i\mathbf{v})$ gives $T\mathbf{u}=a\mathbf{u}-b\mathbf{v}$ and $T\mathbf{v}=b\mathbf{u}+a\mathbf{v}$.\\
Thus $\mathcal{U}=\operatorname{span}(\mathbf{u},\mathbf{v})$ is nonzero, has dimension $1$ or $2$, and is invariant under $T$.\\
\hlt{Structural Commentary}: A complex eigenvector decomposes into two real directions that are closed under $T$.\\
\hlt{Key Insight}: Nonreal eigenvalues become real invariant planes; real eigenvalues become invariant lines.
\end{proof}
\end{theorem}

\begin{figure}[H]
  \centering
  \begin{tikzpicture}[
  >=Stealth,
  every node/.style={font=\small},
  box/.style={draw=gray!55, rounded corners=3pt, inner xsep=7pt, inner ysep=4pt, align=center},
  arrow/.style={->, line width=0.85pt, blue!70!black}
]
  \node[box] (eig) at (0,0) {$T_{\mathbb{C}}(\mathbf{u}+i\mathbf{v})$\\$=(a+bi)(\mathbf{u}+i\mathbf{v})$};
  \node[box] (parts) at (3.55,0) {$T\mathbf{u}=a\mathbf{u}-b\mathbf{v}$\\$T\mathbf{v}=b\mathbf{u}+a\mathbf{v}$};
  \node[box] (plane) at (6.90,0) {$\mathcal{U}=\operatorname{span}(\mathbf{u},\mathbf{v})$\\$T\mathcal{U}\subseteq\mathcal{U}$};

  \draw[arrow] (eig) -- (parts);
  \draw[arrow] (parts) -- (plane);
\end{tikzpicture}

  \caption{Separate a complex eigenpair into a real invariant plane.}
  \label{fig:real-invariant-plane-from-complex-eigenpair}
\end{figure}

\begin{perspective} \label{per:real-normal-operators-planes} \hlt{Real Normal Operators Use Planes}\\
Over $\mathbb{R}$, normal operators need not have enough eigenvectors. The replacement for a missing conjugate eigenpair is a two-dimensional invariant plane carrying a rotation-dilation block.
\end{perspective}

\begin{proposition} \label{prop:two-dimensional-real-normal-block} \hlt{Two-Dimensional Real Normal Blocks}\\
Let $\dim(\mathcal{V})=2$ over $\mathbb{R}$ and let $T\in\LL(\mathcal{V})$.\\
Then $T$ is normal but not self-adjoint if and only if, in every orthonormal basis, $[T]=
\begin{bmatrix}
a&-b\\
b&a
\end{bmatrix}$, $b\neq0$.\\
Moreover, some orthonormal basis makes $b>0$.
\begin{proof}
\begin{enumerate}[label=\roman*.]
\setlength{\itemsep}{0pt}
\item[$(\Rightarrow)$] In an orthonormal basis, write $A=[T]=
\begin{bmatrix}
a&c\\
b&d
\end{bmatrix}$.\\
Normality gives $A^tA=AA^t$. Comparing the $(1,1)$ and $(2,2)$ entries gives $b^2=c^2$. If $c=b$, then $A=A^t$, impossible because $T$ is not self-adjoint. Thus $c=-b$ and $b\neq0$. Comparing off-diagonal entries gives $-ab+bd=ab-bd$, so $d=a$. Since the original orthonormal basis was arbitrary, every orthonormal basis has the displayed form.
\item[$(\Leftarrow)$] A matrix $\begin{bmatrix}a&-b\\ b&a\end{bmatrix}$ with $b\neq0$ satisfies $A^tA=AA^t$ and is not symmetric. Hence $T$ is normal but not self-adjoint. If $b<0$, replace the second basis vector by its negative to make $b>0$.
\end{enumerate}
\hlt{Structural Commentary}: In real dimension $2$, non-self-adjoint normality is exactly a scalar plus a rotation block.\\
\hlt{Key Insight}: Normality forces $c=-b$ and $d=a$.
\end{proof}
\end{proposition}

\begin{theorem} \label{thm:real-normal-operator-block-diagonal-form} \hlt{Real Normal Operator Block Form}\\
Let $\mathcal{V}$ be finite-dimensional over $\mathbb{R}$ and let $T\in\LL(\mathcal{V})$. Then $T$ is normal if and only if some orthonormal basis makes $[T]$ block diagonal, with each diagonal block either $1$-by-$1$ or of the form $\begin{bmatrix}
a&-b\\
b&a
\end{bmatrix}$, $b > 0$.
\begin{proof}
\begin{enumerate}[label=\roman*.]
\setlength{\itemsep}{0pt}
\item[$(\Rightarrow)$] Induct on $n=\dim(\mathcal{V})$. The cases $n\leq1$ are immediate.\\
For $n\geq2$, Theorem \ref{thm:real-operator-invariant-line-or-plane} gives an invariant subspace $\mathcal{U}$ with $\dim(\mathcal{U})\in\{1,2\}$.\\
By Corollary \ref{cor:normal-operator-invariant-subspaces-reduce}, $\mathcal{U}^{\perp}$ is invariant under $T$ and both restrictions are normal.\\
If $\dim(\mathcal{U})=1$, take a unit vector in $\mathcal{U}$. If $\dim(\mathcal{U})=2$, use the Real Spectral Theorem \ref{thm:real-spectral-theorem} when $T|_{\mathcal{U}}$ is self-adjoint, giving scalar $1$-by-$1$ blocks on $\mathcal{U}$, and use Proposition \ref{prop:two-dimensional-real-normal-block} otherwise.\\
By induction, choose the required orthonormal block basis on $\mathcal{U}^{\perp}$. Concatenating it with the basis chosen on $\mathcal{U}$ gives an orthonormal basis of $\mathcal{V}$ because $\mathcal{V}=\mathcal{U}\oplus\mathcal{U}^{\perp}$.
\item[$(\Leftarrow)$] Each allowed block is normal. If $[T]$ is the block diagonal sum of these blocks, then $[T]^*[T]$ and $[T][T]^*$ are the corresponding block diagonal sums of $B^*B$ and $BB^*$, hence are equal block by block.
\end{enumerate}
\hlt{Structural Commentary}: Normality lets invariant lines or planes split orthogonally.\\
\hlt{Key Insight}: Real normal operators split into scalar lines and rotation-dilation planes.
\end{proof}
\end{theorem}

\begin{figure}[H]
  \centering
  \begin{tikzpicture}[
  >=Stealth,
  every node/.style={font=\small},
  box/.style={draw=gray!55, rounded corners=3pt, inner xsep=8pt, inner ysep=6pt, align=center},
  arrow/.style={->, line width=0.85pt, blue!70!black}
]
  \node[box] (normal) at (0,0) {normal $T$};
  \node[box] (sum) at (3.6,0) {$\mathcal{V}=\bigoplus L_j\oplus\bigoplus P_k$};
  \node[box] (blocks) at (7.8,0) {$[T]=\operatorname{diag}(\lambda_j,R_k)$\\[-1pt]
  $R_k=\begin{bmatrix}a_k&-b_k\\ b_k&a_k\end{bmatrix}$};

  \draw[arrow] (normal.east) -- (sum.west);
  \draw[arrow] (sum.east) -- (blocks.west);
\end{tikzpicture}

  \caption{Orthogonal line/plane split into scalar and rotation blocks.}
  \label{fig:real-normal-block-structure}
\end{figure}

\begin{example} \label{ex:plane-rotation-isometry} \hlt{Plane Rotations}\\
For $\theta\in\mathbb{R}$, the real matrix $R_{\theta}=
\begin{bmatrix}
\cos\theta&-\sin\theta\\
\sin\theta&\cos\theta
\end{bmatrix}$ is an isometry of $\mathbb{R}^2$. If $\theta$ is not an integer multiple of $\pi$, then $R_{\theta}$ has no real eigenvalue, which is why real isometries require two-dimensional rotation blocks.
\end{example}

\begin{theorem} \label{thm:real-isometry-block-diagonal-form} \hlt{Real Isometry Block Form}\\
Let $\mathcal{V}$ be finite-dimensional over $\mathbb{R}$ and let $S\in\LL(\mathcal{V})$. Then $S$ is an isometry if and only if some orthonormal basis makes $[S]$ block diagonal, with each diagonal block one of $[1]$, $[-1]$, $\begin{bmatrix}
\cos\theta&-\sin\theta\\
\sin\theta&\cos\theta
\end{bmatrix}$ with $0<\theta<\pi$.
\begin{proof}
\begin{enumerate}[label=\roman*.]
\setlength{\itemsep}{0pt}
\item[$(\Rightarrow)$] Since $S$ is an isometry, $S^*S=I_{\mathcal{V}}$ by Theorem \ref{thm:finite-dimensional-isometry-iff-adjoint-product-identity}. Finite-dimensional invertibility gives $SS^*=I_{\mathcal{V}}$, so $S$ is normal. Apply Theorem \ref{thm:real-normal-operator-block-diagonal-form}. A $1$-by-$1$ block $[\lambda]$ has $\lambda^2=1$, so $\lambda=\pm1$.\\
A $2$-by-$2$ block $\begin{bmatrix}a&-b\\b&a\end{bmatrix}$ has $a^2+b^2=1$ and $b>0$, so $a=\cos\theta$ and $b=\sin\theta$ for a unique $\theta\in(0,\pi)$.
\item[$(\Leftarrow)$] Each displayed block $B$ satisfies $B^tB=I$ on its coordinate subspace. Therefore the block diagonal matrix satisfies $[S]^t[S]=I$, so $S$ preserves norms and is an isometry.
\end{enumerate}
\hlt{Structural Commentary}: Isometries are the unit-length special case of real normal blocks.\\
\hlt{Key Insight}: Scalar blocks reduce to $\pm1$, and planar normal blocks reduce to rotations.
\end{proof}
\end{theorem}

\begin{figure}[H]
  \centering
  \begin{tikzpicture}[
  >=Stealth,
  every node/.style={font=\small},
  box/.style={draw=gray!55, rounded corners=3pt, inner xsep=8pt, inner ysep=6pt, align=center},
  arrow/.style={->, line width=0.85pt, blue!70!black}
]
  \node[box] (iso) at (0,0) {isometry $S$};
  \node[box] (blocks) at (3.3,0) {orthogonal sum\\$[1]$, $[-1]$, $R_{\theta}$};
  \node[box] (rot) at (7.2,0) {$R_{\theta}=\begin{bmatrix}\cos\theta&-\sin\theta\\ \sin\theta&\cos\theta\end{bmatrix}$};

  \draw[arrow] (iso) -- (blocks);
  \draw[arrow] (blocks) -- (rot);
\end{tikzpicture}

  \caption{Normal form: fixed/flipped lines and rotation planes.}
  \label{fig:real-isometry-blocks}
\end{figure}

\begin{proposition} \label{prop:real-eigenvalues-complexification} \hlt{Real Eigenvalues of a Complexification}\\
Let $\mathcal{V}$ be real, let $T\in\LL(\mathcal{V})$, and let $\lambda\in\mathbb{R}$. Then $\lambda$ is an eigenvalue of $T$ if and only if it is an eigenvalue of $T_{\mathbb{C}}$.
\begin{proof}
\begin{enumerate}[label=\roman*.]
\setlength{\itemsep}{0pt}
\item[$(\Rightarrow)$] If $T\mathbf{v}=\lambda\mathbf{v}$ with $\mathbf{v}\neq\mathbf{0}$, then $T_{\mathbb{C}}(\mathbf{v}+i\mathbf{0})=\lambda(\mathbf{v}+i\mathbf{0})$.
\item[$(\Leftarrow)$] If $T_{\mathbb{C}}(\mathbf{u}+i\mathbf{v})=\lambda(\mathbf{u}+i\mathbf{v})$ and $\mathbf{u}+i\mathbf{v}\neq\mathbf{0}$, then $T\mathbf{u}=\lambda\mathbf{u}$ and $T\mathbf{v}=\lambda\mathbf{v}$.\\
Since $\mathbf{u}$ or $\mathbf{v}$ is nonzero, $T$ has eigenvalue $\lambda$.
\end{enumerate}
\hlt{Structural Commentary}: Real eigenvalues survive complexification exactly.\\
\hlt{Key Insight}: For real $\lambda$, real and imaginary parts decouple.
\end{proof}
\end{proposition}

\begin{lemma} \label{lem:complexification-conjugate-generalized-eigenspaces} \hlt{Conjugation Transfers Generalized Eigenspaces}\\
Let $\mathcal{V}$ be real, let $T\in\LL(\mathcal{V})$, let $\lambda\in\mathbb{C}$, and let $k\geq0$. Then $(T_{\mathbb{C}}-\lambda I)^k(\mathbf{u}+i\mathbf{v})=\mathbf{0}$ $\Longleftrightarrow$ $(T_{\mathbb{C}}-\overline{\lambda} I)^k(\mathbf{u}-i\mathbf{v})=\mathbf{0}$.
\begin{proof}
Let $\Gamma(\mathbf{u}+i\mathbf{v})=\mathbf{u}-i\mathbf{v}$.\\
Then $\Gamma$ is a conjugate-linear bijection, $\Gamma T_{\mathbb{C}}=T_{\mathbb{C}}\Gamma$, and $\Gamma\bigl((T_{\mathbb{C}}-\lambda I)^k\mathbf{x}\bigr) =(T_{\mathbb{C}}-\overline{\lambda} I)^k\Gamma\mathbf{x}$.
\begin{enumerate}[label=\roman*.]
\setlength{\itemsep}{0pt}
\item[$(\Rightarrow)$] Apply the displayed identity to $\mathbf{x}=\mathbf{u}+i\mathbf{v}$.
\item[$(\Leftarrow)$] Apply the first direction with $\overline{\lambda}$ in place of $\lambda$ and $\mathbf{u}-i\mathbf{v}$ in place of $\mathbf{u}+i\mathbf{v}$.
\end{enumerate}
\hlt{Structural Commentary}: Conjugation commutes with $T_{\mathbb{C}}$ and conjugates scalars.\\
\hlt{Key Insight}: Conjugation pairs generalized eigenspaces.
\end{proof}
\end{lemma}

\begin{corollary} \label{cor:complexification-conjugate-eigenvalue-multiplicities} \hlt{Conjugate Eigenvalues Have Equal Multiplicity}\\
Let $\mathcal{V}\neq\{\mathbf{0}\}$ be finite-dimensional over $\mathbb{R}$ and let $T\in\LL(\mathcal{V})$.\\
If $\lambda$ is an eigenvalue of $T_{\mathbb{C}}$, then $\overline{\lambda}$ is an eigenvalue of $T_{\mathbb{C}}$, and $\dim_{\mathbb{C}}\mathcal{G}_{\lambda}(T_{\mathbb{C}}) =\dim_{\mathbb{C}}\mathcal{G}_{\overline{\lambda}}(T_{\mathbb{C}})$.
\begin{proof}
Lemma \ref{lem:complexification-conjugate-generalized-eigenspaces} with $k=1$ sends a nonzero eigenvector for $\lambda$ to a nonzero eigenvector for $\overline{\lambda}$. Thus $\overline{\lambda}$ is an eigenvalue. With $k=\dim_{\mathbb{C}}\mathcal{V}_{\mathbb{C}}$, Theorem \ref{thm:generalized-eigenspace-stabilized-kernel} identifies generalized eigenspaces with the corresponding stabilized kernels. By Lemma \ref{lem:complexification-conjugate-generalized-eigenspaces}, $\Gamma(\mathbf{u}+i\mathbf{v})=\mathbf{u}-i\mathbf{v}$ restricts to conjugate-linear bijection $\mathcal{G}_{\lambda}(T_{\mathbb{C}}) \longrightarrow \mathcal{G}_{\overline{\lambda}}(T_{\mathbb{C}})$.\\
Conjugate-linear bijections preserve complex dimension, so the multiplicities agree.\\
\hlt{Structural Commentary}: Conjugation mirrors generalized eigenspaces.\\
\hlt{Key Insight}: Nonreal spectral data of a real operator arrives in equal-sized conjugate pairs.
\end{proof}
\end{corollary}

\begin{corollary} \label{cor:odd-dimensional-real-operator-has-eigenvalue} \hlt{Odd-Dimensional Real Operators Have Eigenvalues}\\
Every operator on a nonzero odd-dimensional real vector space has a real eigenvalue.
\begin{proof}
Let $T\in\LL(\mathcal{V})$, where $\dim_{\mathbb{R}}\mathcal{V}$ is odd. By Proposition \ref{prop:complexification-vector-space-basis}, $\dim_{\mathbb{C}}\mathcal{V}_{\mathbb{C}}$ is odd.\\
The sum of the multiplicities of all eigenvalues of $T_{\mathbb{C}}$ equals this odd dimension by Proposition \ref{prop:sum-of-complex-eigenvalue-multiplicities}. Nonreal eigenvalues contribute in conjugate pairs of equal multiplicity by Corollary \ref{cor:complexification-conjugate-eigenvalue-multiplicities}, hence contribute an even total. Therefore $T_{\mathbb{C}}$ has at least one real eigenvalue, which is an eigenvalue of $T$ by Proposition \ref{prop:real-eigenvalues-complexification}.\\
\hlt{Structural Commentary}: Odd dimension forces a real eigenvalue.\\
\hlt{Key Insight}: Parity forces at least one real spectral direction.
\end{proof}
\end{corollary}

\begin{definition} \label{def:real-operator-characteristic-polynomial} \hlt{Characteristic Polynomial of a Real Operator}\\
Let $\mathcal{V}\neq\{\mathbf{0}\}$ be a finite-dimensional real vector space and let $T\in\LL(\mathcal{V})$.\\
The \hlt{characteristic polynomial} of $T$ is $q_T\defeq q_{T_{\mathbb{C}}}$, the characteristic polynomial of its complexification.
\end{definition}

\begin{theorem} \label{thm:real-operator-characteristic-polynomial-properties} \hlt{Characteristic Polynomial over $\mathbb{R}$}\\
Let $\mathcal{V}\neq\{\mathbf{0}\}$ be a finite-dimensional real vector space, let $T\in\LL(\mathcal{V})$, and let $q_T$ be its characteristic polynomial.
\begin{enumerate}[label=(\roman*)]
\setlength{\itemsep}{0pt}
\item $q_T\in\mathbb{R}[z]$.
\item $\deg q_T=\dim_{\mathbb{R}}\mathcal{V}$.
\item The real zeros of $q_T$ are precisely the eigenvalues of $T$.
\item $q_T(T)=0$.
\end{enumerate}
\begin{proof}
\begin{enumerate}[label=(\roman*)]
\setlength{\itemsep}{0pt}
\item Factor $q_{T_{\mathbb{C}}}$ over $\mathbb{C}$. Real eigenvalues give real linear factors. If $\lambda\notin\mathbb{R}$ has multiplicity $m$, then $\overline{\lambda}$ has multiplicity $m$ by Corollary \ref{cor:complexification-conjugate-eigenvalue-multiplicities}, and $(z-\lambda)^m(z-\overline{\lambda})^m=\bigl(z^2-2\operatorname{Re}(\lambda)z+\abs{\lambda}^2\bigr)^m$, which is real.
\item Proposition \ref{prop:operator-characteristic-polynomial-degree-roots} and Proposition \ref{prop:complexification-vector-space-basis} give $\deg q_T=\dim_{\mathbb{R}}\mathcal{V}$.
\item Real zeros of $q_T=q_{T_{\mathbb{C}}}$ are real eigenvalues of $T_{\mathbb{C}}$  (Prop \ref{prop:operator-characteristic-polynomial-degree-roots}), hence eigenvalues of $T$ (Prop \ref{prop:real-eigenvalues-complexification}).
\item Cayley--Hamilton over $\mathbb{C}$ gives $q_T(T_{\mathbb{C}})=0$. By (i) and Proposition \ref{prop:complexification-operator-basic-properties}, $q_T(T_{\mathbb{C}})=(q_T(T))_{\mathbb{C}}$.\\
If $q_T(T)$ were nonzero on some $\mathbf{u}\in\mathcal{V}$, then $(q_T(T))_{\mathbb{C}}(\mathbf{u}+i\mathbf{0})=q_T(T)\mathbf{u}+i\mathbf{0}$ would be nonzero, a contradiction. Hence $q_T(T)=0$.
\end{enumerate}
\hlt{Structural Commentary}: Nonreal factors pair into real quadratics.\\
\hlt{Key Insight}: Complexify, then pair conjugates.
\end{proof}
\end{theorem}

\begin{figure}[H]
  \centering
  \begin{tikzpicture}[
  >=Stealth,
  every node/.style={font=\small},
  box/.style={draw=gray!55, rounded corners=3pt, inner xsep=6.5pt, inner ysep=4.5pt, align=center},
  arrow/.style={->, line width=0.85pt, blue!70!black}
]
  \node[box] (complex) at (0,0) {$(z-\lambda)^m(z-\overline{\lambda})^m$};
  \node[box] (real) at (3.87,0) {$\bigl(z^2-2\operatorname{Re}(\lambda)z+\abs{\lambda}^2\bigr)^m$};
  \node[box] (poly) at (7.17,0) {$q_T\in\mathbb{R}[z]$};

  \draw[arrow] (complex) -- (real);
  \draw[arrow] (real) -- (poly);
\end{tikzpicture}

  \caption{Conjugate factors pair into real quadratics.}
  \label{fig:real-characteristic-conjugate-pairs}
\end{figure}

\begin{corollary} \label{cor:real-characteristic-polynomial-multiple-minimal-polynomial} \hlt{$q_T$ Is a Multiple of $\mu_T$ over $\mathbb{R}$}\\
Let $\mathcal{V}\neq\{\mathbf{0}\}$ be a finite-dimensional real vector space and let $T\in\LL(\mathcal{V})$.\\
Then $\deg\mu_T\leq\dim_{\mathbb{R}}\mathcal{V}$, and $q_T$ is a polynomial multiple of $\mu_T$ in $\mathbb{R}[z]$.
\begin{proof}
Theorem \ref{thm:real-operator-characteristic-polynomial-properties} gives $q_T(T)=0$.\\
By Theorem \ref{thm:operator-annihilator-iff-multiple-minimal-polynomial}, $\mu_T$ divides $q_T$ in $\mathbb{R}[z]$. The degree bound follows from $\deg q_T=\dim_{\mathbb{R}}\mathcal{V}$.\\
\hlt{Structural Commentary}: Cayley--Hamilton makes $q_T$ a real annihilator of $T$.\\
\hlt{Key Insight}: The minimal polynomial divides every real annihilator.
\end{proof}
\end{corollary}
